% Part: first-order-logic
% Chapter: tableaux
% Section: proving-things

\documentclass[../../../include/open-logic-section]{subfiles}

\begin{document}

\iftag{FOL}
      {\olfileid{fol}{tab}{pro}}
      {\olfileid{pl}{tab}{pro}}

\olsection{\usetoken{P}{tableau} ఉదాహరణలు}

\begin{ex}
!!{sentence} $(!A \land !B) \lif !A$ కోసం ఒక సంవృత !!{tableau}ను
కనుగొందాం.

మొదట సంబంధిత పరికల్పనను !!{tableau} అగ్రభాగంలో రాస్తాం.
\begin{oltableau}
  [\sFmla{\False}{(\formula{A} \land \formula{B}) \lif \formula{A}},
    just = \TAss]
\end{oltableau}

ఒకే పరికల్పన ఉంది; కనుక నియమాన్ని వర్తింపజేయగల !!{signed formula}
ఒక్కటే. (ప్రతి !!{signed formula}కు వర్తించగల నియమం గరిష్ఠంగా
ఒక్కటే ఉంటుంది: అది సంబంధిత సంకేతానికి, !!{sentence} యొక్క
!!{main operator}కు గల నియమం.) ఈ సందర్భంలో మనం
$\TRule{\False}{\lif}$ను వర్తింపజేయాలి.
\begin{oltableau}
  [\sFmla{\False}{(\formula{A} \land \formula{B}) \lif \formula{A}},
    checked, just = \TAss
    [\sFmla{\True}{\formula{A} \land \formula{B}},
      just={\TRule{\False}{\lif}[1]}
      [\sFmla{\False}{\formula{A}}, just={\TRule{\False}{\lif}[1]}]
    ]
  ]
\end{oltableau}
ఏ !!{signed formula}sకు వాటి సంబంధిత నియమాలను వర్తింపజేశామో
గుర్తుంచుకోవడానికి, ఆ వాక్యం పక్కన తనిఖీ గుర్తు పెడతాం. అయితే ఒక
నియమాన్ని అన్ని వివృత శాఖలకూ వర్తింపజేసినప్పుడు \emph{మాత్రమే}
తనిఖీ గుర్తు పెట్టండి. ఒక !!a{signed formula}కు సంబంధించిన నియమం
ప్రతి వివృత శాఖలోనూ వర్తించిన తర్వాత, మళ్లీ దాని దగ్గరకు తిరిగి వచ్చి
ఆ నియమాన్ని వర్తింపజేయాల్సిన అవసరం లేదు. ఇక్కడ ఒకే శాఖ ఉంది కనుక
నియమాన్ని ఒక్కసారి మాత్రమే వర్తింపజేయాలి. (తనిఖీ గుర్తులు టాబ్లోలను
నిర్మించడంలో సౌలభ్యం కోసం మాత్రమే; అవి టాబ్లోల వాక్యనిర్మాణంలో
అధికారిక భాగం కావని గమనించండి.)

ఇప్పుడు నియమాన్ని వర్తింపజేయగల కొత్త !!{signed formula} ఒకటి ఉంది:
వరుస~$2$లోని $\sFmla{\True}{!A \land !B}$. దానికి
$\TRule{\True}{\land}$ నియమాన్ని వర్తింపజేస్తే:
\begin{oltableau}
  [\sFmla{\False}{(\formula{A} \land \formula{B}) \lif \formula{A}},
    checked, just = \TAss
    [\sFmla{\True}{\formula{A} \land \formula{B}},
      just={\TRule{\False}{\lif}[1]}, checked
      [\sFmla{\False}{\formula{A}}, just={\TRule{\False}{\lif}[1]}
        [\sFmla{\True}{\formula{A}}, just={\TRule{\True}{\land}[2]}
          [\sFmla{\True}{\formula{B}}, just={\TRule{\True}{\land}[2]}, close
          ]
        ]
      ]
    ]
  ]
\end{oltableau}
ఇప్పుడు శాఖలో $\sFmla{\True}{!A}$ (వరుస~$4$లో),
$\sFmla{\False}{!A}$ (వరుస~$3$లో) రెండూ ఉన్నాయి కనుక శాఖ
సంవృతమైంది. అదే ఒక్క శాఖ కనుక !!{tableau} కూడా సంవృతమైంది.
$(!A \land !B) \lif !A$ కోసం ఒక సంవృత !!{tableau}ను కనుగొన్నాం.
\end{ex}

\begin{ex}
ఇప్పుడు $(\lnot !A \lor !B) \lif (!A \lif !B)$ కోసం ఒక సంవృత
!!{tableau}ను కనుగొందాం.

సంబంధిత పరికల్పనతో మొదలుపెడతాం:
\begin{oltableau}
  [\sFmla{\False}{(\lnot \formula{A} \lor \formula{B}) \lif
      (\formula{A} \lif \formula{B})}, just=\TAss]
\end{oltableau}
ఈ !!{tableau}లోని ఒక్క !!{signed formula}కు !!{main operator}~$\lif$,
సంకేతం~$\False$ ఉన్నాయి. కనుక దానికి $\TRule{\False}{\lif}$ను
వర్తింపజేసి కిందిది పొందుతాం:
\begin{oltableau}
  [\sFmla{\False}{(\lnot \formula{A} \lor \formula{B})
      \lif (\formula{A} \lif \formula{B})}, just=\TAss, checked
    [\sFmla{\True}{\lnot \formula{A} \lor \formula{B}},
      just={\TRule{\False}{\lif}[1]}
      [\sFmla{\False}{(\formula{A} \lif \formula{B})},
        just={\TRule{\False}{\lif}[1]}
      ]
    ]
  ]
\end{oltableau}
ఇప్పుడు వరుస~$2$కు $\TRule{\True}{\lor}$ను వర్తింపజేయాలా, లేక
వరుస~$3$కు $\TRule{\False}{\lif}$ను వర్తింపజేయాలా అనే ఎంపిక ఉంది.
ప్రతి !!{signed formula}కు సంబంధించిన నియమాన్ని ప్రతి శాఖలోనూ
వర్తింపజేస్తే చాలు; మనం ఎంచుకునే క్రమం ముఖ్యం కాదు. కాబట్టి మొదటి
ఎంపికను తీసుకుందాం. $\TRule{\True}{\lor}$ నియమం !!{tableau}ను
శాఖలుగా విడగొడుతుంది; నియమపు రెండు నిష్కర్షలు రెండు కొత్త శాఖలకు
చేర్చే కొత్త !!{signed formula}s అవుతాయి. దాని ఫలితం:
\begin{oltableau}
  [\sFmla{\False}{(\lnot \formula{A} \lor \formula{B}) \lif
      (\formula{A} \lif \formula{B})}, just=\TAss, checked
    [\sFmla{\True}{\lnot \formula{A} \lor \formula{B}},
      just={\TRule{\False}{\lif}[1]}, checked
      [\sFmla{\False}{(\formula{A} \lif \formula{B})},
        just={\TRule{\False}{\lif}[1]}
        [\sFmla{\True}{\lnot \formula{A}}, just={\TRule{\True}{\lor}[2]}]
        [\sFmla{\True}{\formula{B}}, just={\TRule{\True}{\lor}[2]}]
      ]
    ]
  ]
\end{oltableau}
వరుస~$3$కు $\TRule{\False}{\lif}$ నియమాన్ని ఇంకా వర్తింపజేయలేదు;
ఇప్పుడు చేద్దాం. సమయం ఆదా చేయడానికి రెండు శాఖలకూ ఒకేసారి
వర్తింపజేస్తాం. సంబంధిత నియమాన్ని ప్రతి వివృత శాఖలోనూ
వర్తింపజేసినప్పుడు మాత్రమే ఒక !!a{signed formula} పక్కన తనిఖీ గుర్తు
పెడతామని గుర్తుంచుకోండి. అందువల్ల ఒక నియమం వర్తించే
!!{signed formula}ను కలిగిన ప్రతి శాఖ చివరలోనూ ఆ నియమాన్ని వర్తింపజేయడం మంచి
పద్ధతి. అప్పుడు వివిధ శాఖల్లో కిందికి వెళ్లిన తర్వాత మళ్లీ ఆ
!!{signed formula} దగ్గరకు తిరిగి రావాల్సిన అవసరం ఉండదు.
\begin{oltableau}
  [\sFmla{\False}{(\lnot \formula{A} \lor \formula{B}) \lif
      (\formula{A} \lif \formula{B})}, just=\TAss, checked
    [\sFmla{\True}{\lnot \formula{A} \lor \formula{B}},
      just={\TRule{\False}{\lif}[1]}, checked
      [\sFmla{\False}{(\formula{A} \lif \formula{B})},
        just={\TRule{\False}{\lif}[1]}, checked
        [\sFmla{\True}{\lnot \formula{A}}, just={\TRule{\True}{\lor}[2]}
          [\sFmla{\True}{\formula{A}}, just={\TRule{\False}{\lif}[3]}
            [\sFmla{\False}{\formula{B}}, just={\TRule{\False}{\lif}[3]}]
          ]
        ]
        [\sFmla{\True}{\formula{B}}, just={\TRule{\True}{\lor}[2]}
          [\sFmla{\True}{\formula{A}}, just={\TRule{\False}{\lif}[3]}
            [\sFmla{\False}{\formula{B}}, just={\TRule{\False}{\lif}[3]}, close]
          ]
        ]
      ]
    ]
  ]
\end{oltableau}
కుడి శాఖ ఇప్పుడు సంవృతమైంది. ఎడమ శాఖలో వరుస~$4$కు
$\TRule{\True}{\lnot}$ నియమాన్ని ఇంకా వర్తింపజేయవచ్చు. దాని వల్ల
$\sFmla{\False}{!A}$ వచ్చి ఎడమ శాఖ కూడా సంవృతమవుతుంది:
\begin{oltableau}
  [\sFmla{\False}{(\lnot \formula{A} \lor \formula{B}) \lif
      (\formula{A} \lif \formula{B})}, just=\TAss, checked
    [\sFmla{\True}{\lnot \formula{A} \lor \formula{B}},
      just={\TRule{\False}{\lif}[1]}, checked
      [\sFmla{\False}{(\formula{A} \lif \formula{B})},
        just={\TRule{\False}{\lif}[1]}, checked
        [\sFmla{\True}{\lnot \formula{A}}, just={\TRule{\True}{\lor}[2]}
          [\sFmla{\True}{\formula{A}}, just={\TRule{\False}{\lif}[3]}
            [\sFmla{\False}{\formula{B}}, just={\TRule{\False}{\lif}[3]}
              [\sFmla{\False}{\formula{A}},
                just={\TRule{\True}{\lnot}[4]}, close]
            ]
          ]
        ]
        [\sFmla{\True}{\formula{B}}, just={\TRule{\True}{\lor}[2]}
          [\sFmla{\True}{\formula{A}}, just={\TRule{\False}{\lif}[3]}
            [\sFmla{\False}{\formula{B}},
              just={\TRule{\False}{\lif}[3]}, close]
          ]
        ]
      ]
    ]
  ]
\end{oltableau}
\end{ex}

\begin{ex}
ఎన్ని !!{signed formula}sనైనా పరికల్పనలుగా తీసుకుని !!{tableau}sను
ఇవ్వవచ్చు. శాఖలుగా విడగొట్టే ఒకటి కంటే ఎక్కువ నియమాలను
వర్తింపజేయడం కూడా తరచుగా అవసరం; సాధారణంగా ఒక !!a{tableau}కు ఎన్ని
శాఖలైనా ఉండవచ్చు. ఉదాహరణకు,
$\{\sFmla{\True}{!A \lor (!B \land !C)}, \sFmla{\False}{(!A \lor !B) \land (!A
  \lor !C)}\}$ కోసం ఒక !!a{tableau}ను పరిశీలించండి. మొదటి పరికల్పనకు
$\TRule{\True}{\lor}$ను వర్తింపజేసి మొదలుపెడతాం:
\begin{oltableau}
  [\sFmla{\True}{\formula{A} \lor (\formula{B} \land \formula{C})},
    just=\TAss, checked
    [\sFmla{\False}{(\formula{A} \lor \formula{B}) \land
        (\formula{A} \lor \formula{C})}, just=\TAss
      [\sFmla{\True}{\formula{A}}, just={\TRule{\True}{\lor}[1]}]
      [\sFmla{\True}{\formula{B} \land \formula{C}},
        just={\TRule{\True}{\lor}[1]}]
    ]
  ]
\end{oltableau}
ఇప్పుడు వరుస~$2$కు $\TRule{\False}{\land}$ నియమాన్ని
వర్తింపజేయవచ్చు. రెండు శాఖలకూ ఒకేసారి వర్తింపజేస్తాం; అందువల్ల
వరుస~$2$కు తనిఖీ గుర్తు పెట్టవచ్చు:
\begin{oltableau}
  [\sFmla{\True}{\formula{A} \lor (\formula{B} \land \formula{C})},
    just=\TAss, checked
    [\sFmla{\False}{(\formula{A} \lor \formula{B}) \land
        (\formula{A} \lor \formula{C})}, just=\TAss, checked
      [\sFmla{\True}{\formula{A}}, just={\TRule{\True}{\lor}[1]}
        [\sFmla{\False}{\formula{A} \lor \formula{B}},
          just={\TRule{\False}{\land}[2]}]
        [\sFmla{\False}{\formula{A} \lor \formula{C}},
          just={\TRule{\False}{\land}[2]}]
      ]
      [\sFmla{\True}{\formula{B} \land \formula{C}},
        just={\TRule{\True}{\lor}[1]}
        [\sFmla{\False}{\formula{A} \lor \formula{B}},
          just={\TRule{\False}{\land}[2]}]
        [\sFmla{\False}{\formula{A} \lor \formula{C}},
          just={\TRule{\False}{\land}[2]}]
      ]
    ]
  ]
\end{oltableau}
ఇప్పుడు $!A \lor !B$ను కలిగి ఉన్న అన్ని శాఖలకూ
$\TRule{\False}{\lor}$ను వర్తింపజేయవచ్చు:
\begin{oltableau}
  [\sFmla{\True}{\formula{A} \lor (\formula{B} \land \formula{C})},
    just=\TAss, checked
    [\sFmla{\False}{(\formula{A} \lor \formula{B}) \land
        (\formula{A} \lor \formula{C})}, just=\TAss, checked
      [\sFmla{\True}{\formula{A}}, just={\TRule{\True}{\lor}[1]}
        [\sFmla{\False}{\formula{A} \lor \formula{B}},
          just={\TRule{\False}{\land}[2]}, checked
          [\sFmla{\False}{\formula{A}}, just={\TRule{\False}{\lor}[4]}
            [\sFmla{\False}{\formula{B}},
              just={\TRule{\False}{\lor}[4]}, close]
          ]
        ]
        [\sFmla{\False}{\formula{A} \lor \formula{C}},
          just={\TRule{\False}{\land}[2]}]
      ]
      [\sFmla{\True}{\formula{B} \land \formula{C}},
        just={\TRule{\True}{\lor}[1]}
        [\sFmla{\False}{\formula{A} \lor \formula{B}},
          just={\TRule{\False}{\land}[2]}, checked
          [\sFmla{\False}{\formula{A}}, just={\TRule{\False}{\lor}[4]}
            [\sFmla{\False}{\formula{B}},
              just={\TRule{\False}{\lor}[4]}]
          ]
        ]
        [\sFmla{\False}{\formula{A} \lor \formula{C}},
          just={\TRule{\False}{\land}[2]}]
      ]
    ]
  ]
\end{oltableau}
అన్నిటికంటే ఎడమ శాఖ ఇప్పుడు సంవృతమైంది. ఇప్పుడు $!A \lor !C$కు
$\TRule{\False}{\lor}$ను వర్తింపజేద్దాం:
\begin{oltableau}
  [\sFmla{\True}{\formula{A} \lor (\formula{B} \land \formula{C})},
    just=\TAss, checked
    [\sFmla{\False}{(\formula{A} \lor \formula{B}) \land
        (\formula{A} \lor \formula{C})}, just=\TAss, checked
      [\sFmla{\True}{\formula{A}}, just={\TRule{\True}{\lor}[1]}
        [\sFmla{\False}{\formula{A} \lor \formula{B}},
          just={\TRule{\False}{\land}[2]}, checked
          [\sFmla{\False}{\formula{A}},
            just={\TRule{\False}{\lor}[4]}
            [\sFmla{\False}{\formula{B}},
              just={\TRule{\False}{\lor}[4]}, close]
          ]
        ]
        [\sFmla{\False}{\formula{A} \lor \formula{C}},
          just={\TRule{\False}{\land}[2]}, checked
          [\sFmla{\False}{\formula{A}},
            just={\TRule{\False}{\lor}[4]}, move by=2
            [\sFmla{\False}{\formula{C}},
              just={\TRule{\False}{\lor}[4]}, close]
          ]
        ]
      ]
      [\sFmla{\True}{\formula{B} \land \formula{C}},
        just={\TRule{\True}{\lor}[1]}
        [\sFmla{\False}{\formula{A} \lor \formula{B}},
          just={\TRule{\False}{\land}[2]}, checked
          [\sFmla{\False}{\formula{A}},
            just={\TRule{\False}{\lor}[4]}
            [\sFmla{\False}{\formula{B}},
              just={\TRule{\False}{\lor}[4]}
            ]
          ]
        ]
        [\sFmla{\False}{\formula{A} \lor \formula{C}},
          just={\TRule{\False}{\land}[2]}, checked
          [\sFmla{\False}{\formula{A}},
            just={\TRule{\False}{\lor}[4]}, move by=2
            [\sFmla{\False}{\formula{C}},
              just={\TRule{\False}{\lor}[4]}]
          ]
        ]
      ]
    ]
  ]
\end{oltableau}
స్పష్టత కోసం $\TRule{\False}{\lor}$ను రెండోసారి వర్తింపజేసిన
ఫలితాన్ని కిందికి కదిలించామని గమనించండి. ఇక్కడ సమర్థనలు ఒకేలా
ఉంటాయి కనుక అది అవసరమయ్యేది కాదు.

రెండు శాఖలు ఇంకా వివృతంగా ఉన్నాయి; వరుస~$3$లోని
$\sFmla{\True}{!B \land !C}$కు తనిఖీ గుర్తు ఇంకా లేదు. దానికి
$\TRule{\True}{\land}$ను వర్తింపజేసి సంవృత !!{tableau}ను పొందుతాం:
\begin{oltableau}
  [\sFmla{\True}{\formula{A} \lor (\formula{B} \land \formula{C})},
    just=\TAss, checked
    [\sFmla{\False}{(\formula{A} \lor \formula{B}) \land
        (\formula{A} \lor \formula{C})}, just=\TAss, checked
      [\sFmla{\True}{\formula{A}}, just={\TRule{\True}{\lor}[1]}
        [\sFmla{\False}{\formula{A} \lor \formula{B}},
          just={\TRule{\False}{\land}[2]}, checked
          [\sFmla{\False}{\formula{A}},
            just={\TRule{\False}{\lor}[4]}
            [\sFmla{\False}{\formula{B}},
              just={\TRule{\False}{\lor}[4]}, close]
          ]
        ]
        [\sFmla{\False}{\formula{A} \lor \formula{C}},
          just={\TRule{\False}{\land}[2]}, checked
          [\sFmla{\False}{\formula{A}},
            just={\TRule{\False}{\lor}[4]}
            [\sFmla{\False}{\formula{C}},
              just={\TRule{\False}{\lor}[4]}, close]
          ]
        ]
      ]
      [\sFmla{\True}{\formula{B} \land \formula{C}},
        just={\TRule{\True}{\lor}[1]}, checked
        [\sFmla{\False}{\formula{A} \lor \formula{B}},
          just={\TRule{\False}{\land}[2]}, checked
          [\sFmla{\False}{\formula{A}},
            just={\TRule{\False}{\lor}[4]}
            [\sFmla{\False}{\formula{B}},
              just={\TRule{\False}{\lor}[4]}
              [\sFmla{\True}{\formula{B}},
                just={\TRule{\True}{\land}[3]}
                [\sFmla{\True}{\formula{C}},
                  just={\TRule{\True}{\land}[3]},close
                ]
              ]
            ]
          ]
        ]
        [\sFmla{\False}{\formula{A} \lor \formula{C}},
          just={\TRule{\False}{\land}[2]}, checked
          [\sFmla{\False}{\formula{A}},
            just={\TRule{\False}{\lor}[4]}
            [\sFmla{\False}{\formula{C}},
              just={\TRule{\False}{\lor}[4]}
              [\sFmla{\True}{\formula{B}},
                just={\TRule{\True}{\land}[3]}
                [\sFmla{\True}{\formula{C}},
                  just={\TRule{\True}{\land}[3]},close
                ]
              ]
            ]
          ]
        ]
      ]
    ]
  ]
\end{oltableau}

పోలిక కోసం, నియమాలను వేరే క్రమంలో వర్తింపజేసిన అదే పరికల్పనల
సమితికి సంవృత !!{tableau} ఇది:
\begin{oltableau}
  [\sFmla{\True}{\formula{A} \lor (\formula{B} \land \formula{C})},
    just=\TAss, checked
    [\sFmla{\False}{(\formula{A} \lor \formula{B}) \land
        (\formula{A} \lor \formula{C})}, just=\TAss, checked
      [\sFmla{\False}{\formula{A} \lor \formula{B}},
        just={\TRule{\False}{\land}[2]}, checked
        [\sFmla{\False}{\formula{A}},
          just={\TRule{\False}{\lor}[3]}
          [\sFmla{\False}{\formula{B}},
            just={\TRule{\False}{\lor}[3]}
            [\sFmla{\True}{\formula{A}},
              just={\TRule{\True}{\lor}[1]},close
            ]
            [\sFmla{\True}{\formula{B} \land \formula{C}},
              just={\TRule{\True}{\lor}[1]}, checked
              [\sFmla{\True}{\formula{B}},
                just={\TRule{\True}{\land}[6]}
                [\sFmla{\True}{\formula{C}},
                  just={\TRule{\True}{\land}[6]},close
                ]
              ]
            ]
          ]
        ]
      ]
      [\sFmla{\False}{\formula{A} \lor \formula{C}},
        just={\TRule{\False}{\land}[2]}, checked
        [\sFmla{\False}{\formula{A}},
          just={\TRule{\False}{\lor}[3]}
          [\sFmla{\False}{\formula{C}},
            just={\TRule{\False}{\lor}[3]}
            [\sFmla{\True}{\formula{A}},
              just={\TRule{\True}{\lor}[1]},close
            ]
            [\sFmla{\True}{\formula{B} \land \formula{C}},
              just={\TRule{\True}{\lor}[1]}, checked
              [\sFmla{\True}{\formula{B}},
                just={\TRule{\True}{\land}[6]}
                [\sFmla{\True}{\formula{C}},
                  just={\TRule{\True}{\land}[6]},close
                ]
              ]
            ]
          ]
        ]
      ]
    ]
  ]
\end{oltableau}
\end{ex}

\begin{prob}
కిందివాటికి సంవృత !!{tableau}s ఇవ్వండి:
\begin{enumerate}
\item $\sFmla{\True}{!A \land (!B \land !C)}, \sFmla{\False}{(!A \land !B) \land !C}$.
\item $\sFmla{\True}{!A \lor (!B \lor !C)}, \sFmla{\False}{(!A \lor !B) \lor !C}$.
\item $\sFmla{\True}{!A \lif (!B \lif !C)}, \sFmla{\False}{!B \lif (!A \lif !C)}$.
\item $\sFmla{\True}{!A}, \sFmla{\False}{\lnot\lnot !A}$.
\end{enumerate}
\end{prob}

\begin{prob}
కిందివాటికి సంవృత !!{tableau}s ఇవ్వండి:
\begin{enumerate}
\item $\sFmla{\True}{(!A \lor !B) \lif !C}, \sFmla{\False}{!A \lif !C}$.
\item $\sFmla{\True}{(!A \lif !C) \land (!B \lif !C)}, \sFmla{\False}{(!A \lor !B) \lif !C}$.
\item $\sFmla{\False}{\lnot(!A \land \lnot !A)}$.
\item $\sFmla{\True}{!B \lif !A}, \sFmla{\False}{\lnot !A \lif \lnot !B}$.
\item $\sFmla{\False}{(!A \lif \lnot !A) \lif \lnot !A}$.
\item $\sFmla{\False}{\lnot(!A \lif !B) \lif \lnot !B}$.
\item $\sFmla{\True}{!A \lif !C}, \sFmla{\False}{\lnot (!A \land \lnot !C)}$.
\item $\sFmla{\True}{!A \land \lnot !C}, \sFmla{\False}{\lnot (!A \lif !C)}$.
\item $\sFmla{\True}{!A \lor !B}, \sFmla{\True}{\lnot !B}, \sFmla{\False}{!A}$.
\item $\sFmla{\True}{\lnot !A \lor \lnot !B}, \sFmla{\False}{\lnot(!A \land !B)}$.
\item $\sFmla{\False}{(\lnot !A \land \lnot !B) \lif\lnot(!A \lor !B)}$.
\item $\sFmla{\False}{\lnot(!A \lor !B) \lif (\lnot !A \land \lnot !B)}$.
\end{enumerate}
\sourcecorrection{OLTETAB-002}{తొమ్మిదవ అంశంలో ఆంగ్ల మూలం
$!A \lor !B$, $\lnot !B$లను ఒకే $\True$-చిహ్నిత సూత్రంలో కామాతో
కలిపింది. ఇది చిహ్నిత సూత్ర నిర్వచనానికి సరిపోదు. అవసరమైన మూడు
పరికల్పనలను విడిగా $\sFmla{\True}{!A \lor !B}$,
$\sFmla{\True}{\lnot !B}$, $\sFmla{\False}{!A}$గా రాశాం.}
\end{prob}

\begin{prob}
కిందివాటికి సంవృత !!{tableau}s ఇవ్వండి:
\begin{enumerate}
\item $\sFmla{\True}{\lnot(!A \lif !B)}, \sFmla{\False}{!A}$.
\item $\sFmla{\True}{\lnot(!A \land !B)}, \sFmla{\False}{\lnot !A \lor \lnot !B}$.
\item $\sFmla{\True}{!A \lif !B}, \sFmla{\False}{\lnot !A \lor !B}$.
\item $\sFmla{\False}{\lnot \lnot !A \lif !A}$.
\item $\sFmla{\True}{!A \lif !B}, \sFmla{\True}{\lnot !A \lif !B}, \sFmla{\False}{!B}$.
\item $\sFmla{\True}{(!A \land !B) \lif !C}, \sFmla{\False}{(!A \lif !C) \lor (!B \lif !C)}$.
\item $\sFmla{\True}{(!A \lif !B) \lif !A}, \sFmla{\False}{!A}$.
\item $\sFmla{\False}{(!A \lif !B) \lor (!B \lif !C)}$.
\end{enumerate}
\end{prob}

\end{document}
